\chapter{Applications of  Production Theory}


\section{1 Structure of Technology}

Empirical work usually imposes certain restrictions on the structure
of production technology. We consider:

\begin{itemize}
    \item Homotheticity
    \item Separability
    \item Nonjointness
\end{itemize}

Elasticity of substitution, which represents production structure,
is also discussed.

\subsection*{1.1 Input Homotheticity}

\[
V(y)=h(y)\,V(1),
\qquad
h:\mathbb{R}_+^M \to \mathbb{R}_+.
\]

Here, $h(y)$ is a scalar and $V(1)\subset \mathbb{R}_+^N$
is a reference input set independent of output.

Under input homotheticity, isoquants are parallel.

\subsection*{Input Distance Function}

\[
D^I(x,y)=\frac{D^I(x)}{h(y)}.
\]

\subsection*{Proof}

\[
D^I(x,y)
=
\max \{\lambda : x \in \lambda V(y)\}
\]
\[
=
\max \{\lambda : x \in \lambda h(y)V(1)\}
\]
\[
=
\max \left\{ \frac{\lambda'}{h(y)} : x \in \lambda' V(1) \right\}
\]
\[
=
\frac{1}{h(y)} \max \{\lambda' : x \in \lambda' V(1)\}
=
\frac{D^I(x)}{h(y)}.
\quad \blacksquare
\]

\subsection*{Cost Function}

\[
c(w,y)=h(y)c(w).
\]

\subsection*{Proof}

\[
c(w,y)
=
\min \{w\cdot x : D^I(x,y)\ge 1\}
\]
\[
=
\min \left\{w\cdot x : \frac{D^I(x)}{h(y)} \ge 1 \right\}
=
\min \{w\cdot x : D^I(x)\ge h(y)\}.
\]

Using homogeneity of degree 1,
\[
=
h(y)\min \{w\cdot x' : D^I(x')\ge 1\}
=
h(y)c(w).
\quad \blacksquare
\]

\subsection*{Profit Function}

\[
\Pi(p,w)
=
\max_y \{p\cdot y - c(w,y)\}
=
\max_y \{p\cdot y - h(y)c(w)\}.
\]

\[
=
c(w)\max_y \left\{ \frac{p}{c(w)}\cdot y - h(y) \right\}
=
c(w)\,\Pi\!\left(\frac{p}{c(w)}\right).
\quad \blacksquare
\]



\subsection*{1.2 Input Separability}

Under input separability, it is possible to construct a single aggregate input.

The production process can be viewed as a two-step procedure:
\begin{itemize}
    \item First, inputs $x$ are aggregated into a scalar index $g(x)$,
    \item Second, output is determined from this aggregate input.
\end{itemize}

\[
y \in Y(x) \quad \Longleftrightarrow \quad y \in Y(g(x)).
\]

\[
g(x): \mathbb{R}_+^N \to \mathbb{R}_+,
\qquad
y: \mathbb{R}_+ \to \mathbb{R}_+^M.
\]

\subsection*{Output Distance Function}

\[
D^O(x,y)=D^O(g(x),y).
\]

\subsection*{Proof}

\[
D^O(x,y)
=
\min \{\theta : y \in \theta Y(x)\}
\]
\[
=
\min \{\theta : y \in \theta Y(g(x))\}
=
D^O(g(x),y).
\quad \blacksquare
\]

\subsection*{Revenue Function}

\[
R(p,x)
=
\max \{p\cdot y : y \in Y(g(x))\}
=
R(p,g(x)).
\]

\subsection*{Profit Function}

\[
\Pi(p,w)
=
\max_g \{R(p,g) - c(w,g)\},
\]

where $g$ is the aggregate input and $c(w,g)$ is the cost required to produce $g$.

\medskip

\textbf{Implication:}

Under input separability, the marginal rates of technical substitution (MRTSs)
are independent of output $y$.



\subsection*{(continued)}

\subsection*{Implication for MRTS}

\[
D^O(x,y)=D^O(g(x),y).
\]

Differentiating,
\[
dD^O(x,y)
=
dD^O(g(x),y)
=
\frac{\partial D^O}{\partial g} \sum_{i=1}^N \frac{\partial g(x)}{\partial x_i} dx_i.
\]

Thus,
\[
\frac{\partial D^O(x,y)}{\partial x_i}
=
\frac{\partial D^O(g(x),y)}{\partial g}
\cdot
\frac{\partial g(x)}{\partial x_i}.
\]

Hence the marginal rate of technical substitution is
\[
\text{MRTS}_{ij}
=
-
\frac{\partial D^O/\partial x_i}{\partial D^O/\partial x_j}
=
-
\frac{\frac{\partial g}{\partial x_i}}{\frac{\partial g}{\partial x_j}}.
\]

\medskip

\textbf{Result:}

MRTS is independent of $y$.

\medskip

Very often, separability holds only for a subset of inputs:
\[
D^O(x,y)=D^O(g_1(x_1),x_2,y).
\]

\subsection*{1.3 Input Nonjointness}

When multiple outputs are produced, production technology exhibits
input nonjointness if each output can be produced independently
and the collection of single-output technologies fully represents
the multi-output technology (Lau, 1972).

\[
V(y)=\sum_{i=1}^M V_i(y_i),
\]

where
\[
V_i(y_i)=\{x_i : y_i \le f_i(x_i)\}.
\]

\subsection*{Input Distance Function}

\[
D^I(x,y)
=
\max \left\{
\min_{i=1,\dots,M} D_i(x_i,y_i)
\ :\ \sum_{i=1}^M x_i = x
\right\},
\]

where
\[
D_i(x_i,y_i)=\max \{\lambda : x_i \in \lambda V_i(y_i)\}.
\]

\subsection*{Proof}

\[
D^I(x,y)
=
\max \{\lambda : x \in \lambda V(y)\}
\]
\[
=
\max \left\{\lambda : x \in \lambda \sum_{i=1}^M V_i(y_i)\right\}
\]
\[
=
\max \left\{
\lambda :
x_i \in \lambda V_i(y_i),\ i=1,\dots,M,\ \sum_{i=1}^M x_i = x
\right\}
\]
\[
=
\max \left\{
\min_{i=1,\dots,M} D_i(x_i,y_i)
\ :\ \sum_{i=1}^M x_i = x
\right\}.
\quad \blacksquare
\]

\subsection*{Cost Function}

\[
c(w,y)
=
\sum_{i=1}^M c_i(w,y_i),
\]

i.e., the sum of single-output cost functions (no economies of scope).

\subsection*{Proof}

\[
c(w,y)
=
\min \{w\cdot x : x \in V(y)\}
\]
\[
=
\min \left\{w\cdot x : x \in \sum_{i=1}^M V_i(y_i)\right\}
\]
\[
=
\min \left\{
\sum_{i=1}^M w\cdot x_i
:\ x_i \in V_i(y_i)
\right\}
\]
\[
=
\sum_{i=1}^M \min \{w\cdot x_i : x_i \in V_i(y_i)\}
=
\sum_{i=1}^M c_i(w,y_i).
\quad \blacksquare
\]

\subsection*{Profit Function}

\[
\Pi(p,w)
=
\sum_{i=1}^M \Pi_i(p_i,w),
\]

i.e., the sum of single-output profit functions.


\subsection*{(continued)}

\subsection*{Profit Function under Nonjointness}

\[
\Pi(p,w)
=
\max_y \left\{ p\cdot y - \sum_{i=1}^M c_i(w,y_i) \right\}
\]
\[
=
\sum_{i=1}^M \max_{y_i} \{ p_i y_i - c_i(w,y_i) \}
=
\sum_{i=1}^M \Pi_i(p_i,w).
\]

\medskip

\textbf{Implication (Test for Nonjointness):}

Under input nonjointness,
\[
\frac{\partial y_i(p,w)}{\partial p_j}
=
\frac{\partial^2 \Pi(p,w)}{\partial p_i \partial p_j}
=
0
\quad \text{for } i \ne j.
\]

\medskip

That is, there is \textbf{no cross-price effect}.

This provides an empirical test for nonjointness.

\section{1.4 Elasticity of Substitution}

\subsection*{1.4.1 Definition}

Consider a two-input production function:
\[
y = f(x_1, x_2).
\]

The elasticity of substitution provides a unit-free measure of the
substitutability between inputs.

\[
\sigma
=
\frac{d(x_2/x_1)}{d(\text{MRTS})}
\cdot
\frac{\text{MRTS}}{x_2/x_1}.
\]

Equivalently,
\[
\sigma
=
\frac{d \ln(x_2/x_1)}{d \ln(\text{MRTS})}.
\]

Since
\[
\text{MRTS}
=
\frac{f_1}{f_2},
\]
we can write
\[
\sigma
=
\frac{d \ln(x_2/x_1)}{d \ln(f_1/f_2)}.
\]

\medskip

\textbf{Interpretation:}

$\sigma$ measures the responsiveness of the input ratio to changes in MRTS.

\medskip

Suppose the optimal choice moves from point $A$ to point $B$.
The slopes at these points represent the MRTSs.

Then $\sigma$ captures how much the input ratio changes relative to
the change in MRTS.



\subsection*{(continued)}

Suppose the choice moves from point $A$ to point $B$ on an isoquant.
Then
\[
\frac{d(x_2/x_1)}{d(f_1/f_2)}
\]
is given by the ratio of the proportional change in the input ratio
to the proportional change in MRTS.

Thus, $\sigma$ measures the curvature of the isoquant.

\medskip

Equivalently, $\sigma$ can be rewritten as
\[
\sigma
=
-\frac{f_1 f_2 (x_1 f_1 + x_2 f_2)}
{x_1 x_2 \left( f_{11} f_2^2 - 2 f_{12} f_1 f_2 + f_{22} f_1^2 \right)}.
\]

\subsection*{Derivation Sketch}

\begin{enumerate}
    \item Note that
    \[
    \frac{d}{dx_2}\left(\frac{x_2}{x_1}\right)
    =
    \frac{x_1 - x_2 \frac{dx_1}{dx_2}}{x_1^2}.
    \]

    Since along an isoquant
    \[
    f_1\,dx_1 + f_2\,dx_2 = 0,
    \]
    we have
    \[
    \frac{dx_1}{dx_2} = -\frac{f_2}{f_1}.
    \]

    \item Similarly,
    \[
    \frac{d}{dx_2}\left(\frac{f_1}{f_2}\right)
    =
    \frac{\left(f_{11}\frac{dx_1}{dx_2} + f_{12}\right)f_2
    -
    f_1\left(f_{21}\frac{dx_1}{dx_2} + f_{22}\right)}
    {f_2^2}.
    \]

    Using $f_{12}=f_{21}$ and substituting
    \[
    \frac{dx_1}{dx_2}=-\frac{f_2}{f_1},
    \]
    this simplifies to the denominator term above.

    \item Combining the two expressions yields the formula for $\sigma$.
\end{enumerate}

\subsection*{Bordered Hessian Form}

Let
\[
F =
\begin{vmatrix}
0 & f_1 & f_2 \\
f_1 & f_{11} & f_{12} \\
f_2 & f_{21} & f_{22}
\end{vmatrix}.
\]

Then $\sigma$ can also be written as
\[
\sigma
=
\frac{x_1 f_1 + x_2 f_2}{x_1 x_2}\cdot \frac{f_1 f_2}{F}.
\]

Since $f_{12}=f_{21}$, the measure $\sigma$ is symmetric in the two inputs.

If the production function is quasiconcave, so that the isoquant is convex
to the origin, then
\[
\sigma > 0.
\]

\subsection*{Cost-Minimization Interpretation}

From the cost-minimization problem,
\[
w_i = \lambda f_i,
\]
so
\[
\frac{f_1}{f_2} = \frac{w_1}{w_2}.
\]

Hence,
\[
\sigma
=
\frac{d\ln(x_2/x_1)}{d\ln(f_1/f_2)}
=
\frac{d\ln(x_2/x_1)}{d\ln(w_1/w_2)}.
\]

Thus, $\sigma$ can also be interpreted as the elasticity of the input ratio
with respect to the ratio of input prices.

\subsection*{More than Two Inputs}

When there are more than two inputs, there are several measures of
substitution elasticity.

The most popular is the \textbf{Allen--Uzawa elasticity of substitution},
defined by
\[
\sigma_{ij}
=
\frac{\sum_{k=1}^N x_k f_k}{x_i x_j}\cdot \frac{F_{ij}}{F},
\]
where $F_{ij}$ is the cofactor associated with $f_{ij}$ in the bordered Hessian

\[
F =
\begin{vmatrix}
0 & f_1 & \cdots & f_N \\
f_1 & f_{11} & \cdots & f_{1N} \\
\vdots & \vdots & \ddots & \vdots \\
f_N & f_{N1} & \cdots & f_{NN}
\end{vmatrix}.
\]

The Allen--Uzawa elasticity is also symmetric:
\[
\sigma_{ij} = \sigma_{ji}.
\]

Goods are \textbf{complements} if $\sigma_{ij}<0$ and
\textbf{substitutes} if $\sigma_{ij}>0$.

It can be shown that at least one $\sigma_{ij}$ must be positive.

You should also be able to derive the relation
\[
\sigma_{ij} = \frac{E_{ij}}{s_j},
\]
where $s_j$ is the cost share of input $x_j$, and $E_{ij}$ is the elasticity
of conditional factor demand.

\subsection*{Example 1}

Consider the production function
\[
y = f(x_1,x_2).
\]

Taking total differentiation of the production function gives
\[
dy - f_1\,dx_1 - f_2\,dx_2 = 0.
\]

From the cost-minimization conditions
\[
w_i = \lambda f_i,
\]
we also obtain
\[
dw_1 - f_1\,d\lambda - \lambda f_{11}\,dx_1 - \lambda f_{12}\,dx_2 = 0,
\]
\[
dw_2 - f_2\,d\lambda - \lambda f_{21}\,dx_1 - \lambda f_{22}\,dx_2 = 0.
\]

In matrix form,
\[
\begin{pmatrix}
0 & \lambda f_1 & \lambda f_2 \\
f_1 & \lambda f_{11} & \lambda f_{12} \\
f_2 & \lambda f_{21} & \lambda f_{22}
\end{pmatrix}
\begin{pmatrix}
d\lambda \\
dx_1 \\
dx_2
\end{pmatrix}
=
\begin{pmatrix}
dy \\
dw_1 \\
dw_2
\end{pmatrix}.
\]


\subsection*{(continued)}

From cost minimization, we have
\[
w_1 x_1 + w_2 x_2 = C(w),
\]
and the elasticity relation can be written as
\[
\frac{\partial x_2}{\partial w_1}
=
- \frac{w_1 x_1 + w_2 x_2}{w_1 x_1} \cdot \frac{x_2}{x_1}
=
- \frac{x_2}{x_1} \cdot \frac{1}{s_1},
\]
where
\[
s_1 = \frac{w_1 x_1}{w_1 x_1 + w_2 x_2}
\]
is the cost share of input $x_1$.

\subsection*{1.4.2 Constant Elasticity of Substitution (CES)}

A CES production function is given by
\[
Q =
\left[
aL^{\rho} + bK^{\rho}
\right]^{\frac{1}{\rho}},
\qquad a + b = 1.
\]

The elasticity of substitution is
\[
\sigma = \frac{1}{1-\rho}.
\]

\subsection*{MRTS}

\[
\text{MRTS}
=
\frac{MPL}{MPK}
=
\frac{a L^{\rho-1}}{b K^{\rho-1}}
=
\frac{a}{b} \left(\frac{L}{K}\right)^{\rho-1}.
\]

Thus,
\[
\frac{K}{L}
=
\left(\frac{b}{a} \cdot \text{MRTS} \right)^{\frac{1}{\rho-1}}.
\]

Taking derivatives,
\[
\sigma
=
\frac{d \ln(K/L)}{d \ln(\text{MRTS})}
=
\frac{1}{1-\rho}.
\]

---

\subsection*{Special Cases}

\begin{enumerate}
    \item \textbf{Perfect Substitutes} ($\rho = 1$)

    \[
    Q = aL + bK.
    \]

    \item \textbf{Cobb--Douglas} ($\rho \to 0$)

    Using l'H\^opital's rule:
    \[
    \ln Q
    =
    \lim_{\rho \to 0}
    \frac{\ln(aL^\rho + bK^\rho)}{\rho}
    \]

    Expanding,
    \[
    \ln Q = a \ln L + b \ln K,
    \]

    hence
    \[
    Q = L^a K^b.
    \]

    \item \textbf{Perfect Complements} ($\rho \to -\infty$)

    Using the generalized mean:
    \[
    \lim_{\rho \to -\infty}
    \left( aL^\rho + bK^\rho \right)^{1/\rho}
    =
    \min(L,K).
    \]
\end{enumerate}

---

\subsection*{Generalized Mean Interpretation}

For $x \in \mathbb{R}_+^N$, define the mean of order $r$:
\[
M_r(x)
=
\left( \sum_{i=1}^N q_i x_i^r \right)^{1/r},
\qquad q_i > 0,\ \sum q_i = 1.
\]

Then:
\[
\lim_{r \to \infty} M_r(x) = \max(x),
\]
\[
\lim_{r \to -\infty} M_r(x) = \min(x).
\]

Thus, CES nests both extremes:
\[
\text{Perfect substitutes} \leftrightarrow \max,
\quad
\text{Perfect complements} \leftrightarrow \min.
\]

---

\subsection*{Remark}

The restriction $a + b = 1$ ensures homogeneity of degree one and
consistent limiting behavior.

\section{2 Parametric Analysis of a Production System}


\section{2 Parametric Analysis of a Production System}

Parametric analysis employs a functional form for production technology.

If the chosen functional form is not sufficiently flexible, it may impose
unintended restrictions on the technology.

\medskip

\textbf{Flexibility} means that the functional form used for empirical
analysis must approximate the true unknown function up to second order.

Let $h^*(z)$ denote the true (unknown) function.
We want to approximate it using a function $h(z)$.

\subsection*{Second-Order Approximation}

We consider three conditions:

\[
(1)\quad h^*(z_0) = h(z_0),
\]
\[
(2)\quad \nabla_z h^*(z_0) = \nabla_z h(z_0),
\]
\[
(3)\quad \nabla^2_z h^*(z_0) = \nabla^2_z h(z_0).
\]

\subsection*{Interpretation}

\begin{itemize}
    \item (1): $h(z)$ approximates the \textbf{level} of $h^*(z)$.
    \item (2): $h(z)$ approximates the \textbf{first derivatives}
    (e.g., input demand or marginal effects).
    \item (3): $h(z)$ approximates the \textbf{second derivatives}
    (curvature properties such as substitution effects).
\end{itemize}

Thus, it is natural to require that any functional form used in empirical
work satisfies (1)--(3).

Functions satisfying these conditions are called
\textbf{second-order differential approximations}.

\subsection*{Number of Independent Effects}

Let $N$ be the dimension of $z$.

\begin{itemize}
    \item Level condition: $1$ effect
    \item First derivatives: $N$ effects
    \item Second derivatives:
    \[
    \frac{1}{2}N(N-1) \quad \text{(off-diagonal)}
    \]
    \[
    +\ N \quad \text{(diagonal)}
    \]
\end{itemize}

Total number of independent effects:
\[
1 + N + \frac{1}{2}N(N-1)
=
\frac{1}{2}(N+1)(N+2).
\]

\subsection*{General Functional Form}

A common flexible functional form (linear in parameters) is:

\[
h(z)
=
\sum_{i=1}^{K} a_i b_i(z),
\]

where:
\begin{itemize}
    \item $b_i(z)$ are known twice continuously differentiable functions,
    \item $a_i$ are parameters to be estimated,
    \item
    \[
    K = \frac{1}{2}(N+1)(N+2).
    \]
\end{itemize}



\subsection*{Matrix Representation of Approximation}

Let $\{b_i(z)\}_{i=1}^K$ be a set of basis functions.
Define the matrix $W$ as:

\[
W =
\begin{pmatrix}
b_1(z) & b_2(z) & \cdots & b_K(z) \\
\frac{\partial b_1(z)}{\partial z_1} & \frac{\partial b_2(z)}{\partial z_1} & \cdots & \frac{\partial b_K(z)}{\partial z_1} \\
\vdots & \vdots & \ddots & \vdots \\
\frac{\partial b_1(z)}{\partial z_N} & \frac{\partial b_2(z)}{\partial z_N} & \cdots & \frac{\partial b_K(z)}{\partial z_N} \\
\frac{\partial^2 b_1(z)}{\partial z_1 \partial z_2} & \frac{\partial^2 b_2(z)}{\partial z_1 \partial z_2} & \cdots & \frac{\partial^2 b_K(z)}{\partial z_1 \partial z_2} \\
\vdots & \vdots & \ddots & \vdots
\end{pmatrix}.
\]

Let the corresponding vector of true function derivatives be:

\[
\begin{pmatrix}
h^*(z) \\
\frac{\partial h^*(z)}{\partial z_1} \\
\vdots \\
\frac{\partial h^*(z)}{\partial z_N} \\
\frac{\partial^2 h^*(z)}{\partial z_1 \partial z_2} \\
\vdots
\end{pmatrix}.
\]

Then $h(z)$ approximates $h^*(z)$ at $z_0$ if:

\begin{itemize}
\item $W(z_0)$ is invertible,
\item $W(z_0)$ has full rank.
\end{itemize}

---

\subsection*{Common Flexible Functional Forms}

\subsubsection*{(1) Generalized Quadratic Function}

\[
h(z)
=
\beta_0
+
\sum_{i=1}^N \beta_i g_i(z_i)
+
\frac{1}{2}
\sum_{i=1}^N \sum_{j=1}^N
\beta_{ij} g_i(z_i) g_j(z_j),
\quad \beta_{ij} = \beta_{ji}.
\]

---

\subsubsection*{(2) Translog Function}

Take a second-order Taylor expansion of $\ln h^*(z)$ around
$z_0 = (1,\dots,1)$:

\[
\ln h(z)
=
\beta_0
+
\sum_{i=1}^N \beta_i \ln z_i
+
\frac{1}{2}
\sum_{i=1}^N \sum_{j=1}^N
\beta_{ij} \ln z_i \ln z_j,
\quad \beta_{ij} = \beta_{ji}.
\]

This is the \textbf{transcendental logarithmic (translog) function}.

---

\subsection*{Cost Function Forms}

\subsubsection*{(1) Generalized Leontief Cost Function}

\[
c(w,y)
=
y \left[
\sum_{i=1}^N \sum_{j=1}^N
\beta_{ij} (w_i w_j)^{1/2}
\right],
\quad \beta_{ij} = \beta_{ji}.
\]

---

\subsubsection*{(2) Translog Cost Function}

\[
\ln c(w,y)
=
\beta_0
+
\sum_{i=1}^N \beta_i \ln w_i
+
\frac{1}{2}
\sum_{i=1}^N \sum_{j=1}^N
\beta_{ij} \ln w_i \ln w_j
+
\beta_y \ln y
+
\frac{1}{2} \beta_{yy} (\ln y)^2
+
\sum_{i=1}^N \beta_{iy} \ln w_i \ln y.
\]

---

\subsection*{Homogeneity Conditions}

For the cost function to be homogeneous of degree 1 in prices:

\[
\sum_{i=1}^N \beta_i = 1,
\]

\[
\sum_{j=1}^N \beta_{ij} = 0 \quad \forall i,
\]

\[
\sum_{i=1}^N \beta_{iy} = 0.
\]


\begin{tcolorbox}

\subsection*{Exercise 1}

Derive the own- and cross-price elasticities of factor demands when the
cost function is either a generalized Leontief function or a translog function.

\end{tcolorbox}

\begin{tcolorbox}

\subsection*{Exercise 2}

Consider a single-output translog cost function.
Show how to impose the restrictions of input homotheticity and input separability,
respectively.

\end{tcolorbox}

\begin{tcolorbox}

\subsection*{Exercise 3}

There are $i=1,\dots,N$ goods, produced with $j=1,\dots,M$ factors.
The production functions are written as
\[
Y_i = f_i(v_i),
\]
where
\[
v_i = (v_{i1},\dots,v_{iM})
\]
is the vector of factor inputs.

\end{tcolorbox}


The production functions are assumed to be positive, increasing, concave,
and homogeneous of degree one in $v_i$ for all $i$.

Let $p_i$ be the price of output $i$ and $w_j$ the price of factor $j$.

The GDP (gross domestic product) function of the economy is defined as
\[
G(p,V)
=
\max_{\{v_i\}}
\sum_{i=1}^N p_i f_i(v_i)
\]
subject to
\[
\sum_{i=1}^N v_i \le V,
\]
where
\[
V=(V_1,\dots,V_M)
\]
is the $(M\times 1)$ endowment vector, and
\[
p=(p_1,\dots,p_N)
\]
is the $(N\times 1)$ output price vector.

\begin{enumerate}
    \item[(i)] Show that $G(p,V)$ is homogeneous of degree one in $p$.
    
    \item[(ii)] Show that $G(p,V)$ is homogeneous of degree one in $V$.
    
    \item[(iii)] Show the following:
    \[
    \frac{\partial G}{\partial p_i} = Y_i,
    \qquad
    \frac{\partial G}{\partial V_j} = w_j,
    \qquad
    \frac{dw_j}{dp_i} = \frac{dY_i}{dV_j}.
    \]
    
    \item[(iv)] Suppose you want to estimate the GDP function using an
    international data set:
    \[
    \ln G
    =
    \alpha_0
    +
    \sum_{i=1}^N \alpha_i \ln p_i
    +
    \sum_{k=1}^M \beta_k \ln V_k
    +
    \frac{1}{2}
    \sum_{i=1}^N \sum_{j=1}^N \gamma_{ij} \ln p_i \ln p_j
    \]
    \[
    \qquad
    +
    \frac{1}{2}
    \sum_{k=1}^M \sum_{l=1}^M \delta_{kl} \ln V_k \ln V_l
    +
    \sum_{i=1}^N \sum_{k=1}^M \phi_{ik} \ln p_i \ln V_k.
    \]
    
    What restrictions must be imposed on $\alpha$, $\beta$, $\gamma$,
    $\delta$, and $\phi$?
    
    \item[(v)] From the GDP function in (iv), derive
    \[
    \frac{\partial Y_i}{\partial V_k}
    \]
    (Rybczynski elasticity).
    
    \item[(vi)] Derive
    \[
    \frac{\partial w_j}{\partial p_i}
    \]
    (Stolper--Samuelson elasticity).
    
    \item[(vii)] Derive
    \[
    \frac{\partial V_l}{\partial V_k}
    \]
    for the cases $k=l$ and $k\ne l$, respectively.
\end{enumerate}